\begin{center}
\begin{minipage}{\linewidth}
\centering
\textbf{同尺度对照：错误轨迹与正确轨迹}\par\smallskip
\begin{tikzpicture}[scale=1.1, every node/.style={font=\small}, line cap=round]
\draw[->,gray] (-4,0)--(4.2,0) node[right] {$x$};
\draw[->,gray] (0,-3.8)--(0,4) node[above] {$y$};
\draw[CStatement,thick] (0,0) ellipse [x radius=3,y radius=2];
\node[below left] at (0,0) {$O$};
\draw[CExplanation,dashed,thick] (0,0) circle [radius={sqrt(13)}];

\draw[CTrap,dash dot,thick] (0,0) circle[radius={sqrt(5)}];
\fill ({sqrt(5)},0) circle(2pt) node[above left] {$F(\sqrt5,0)$};
\node at (4.3,-1.2) {$Q(3,0)$};
\fill (3,0) circle(2pt) ;
\fill (3,2) circle(2pt) node[above right] {$P(3,2)$};
\draw[CProof,thick] (3,-.7)--(3,2.6); \draw[CProof,thick] (-.4,2)--(3.6,2);
\draw (2.75,2)--(2.75,1.75)--(3,1.75);
\node[fill=white] at (0,-4.1) {$OF^2=a^2-b^2=5$};
\node[fill=white] at (0,-4.8) {$OP^2=a^2+b^2=13$};

\end{tikzpicture}
\par 点划圆是误用平方差所得的圆；虚线圆才是交点轨迹。
\end{minipage}
\end{center}
